% Fortsetzung von Kapitel 2 (§ 7). Quelle: vorlesungsmitschrift_gdm1PDF_261114_161111.PDF, PDF-S. 7--13.
\subsection{Das Auswahlaxiom}
\notizquelleNeu{7--13}

\begin{definition}[Auswahlaxiom (AC)]
Das \textbf{Auswahlaxiom} (\emph{Axiom of Choice}, kurz \textbf{AC}) besagt:
Zu jeder nichtleeren Menge $X$ von nichtleeren Mengen existiert eine
\textbf{Auswahlfunktion}, das heißt eine Abbildung
\[
    f:X\longrightarrow\bigcup_{x\in X}x,
    \qquad\text{sodass für alle }x\in X\text{ gilt: }f(x)\in x.
\]
\end{definition}

\begin{beispiel}[Auswahlfunktionen]
\begin{enumerate}
    \item Für $X=\{\{1,2\},\{1\},\{3\},\{1,2,4\}\}$ ist
    $f:X\to\{1,2,3,4\}$ mit
    \[
        \{1,2\}\mapsto1,\quad
        \{1\}\mapsto1,\quad
        \{3\}\mapsto3,\quad
        \{1,2,4\}\mapsto4
    \]
    eine Auswahlfunktion für $X$.
    \item Für die Menge der nichtleeren Teilmengen der natürlichen Zahlen
    $X=\Pot(\N)\setminus\{\emptyset\}$ liefert
    \[
        f:\Pot(\N)\setminus\{\emptyset\}\longrightarrow\N,\qquad
        Z\longmapsto\min(Z)
    \]
    (die Zuordnung zum Minimum der Elemente von $Z$) eine explizite
    Auswahlfunktion. (In der Mitschrift kurz als $\Pot(\N)$ notiert;
    für die leere Menge existiert kein Minimum.)
    \item Für $X=\Pot(\R)\setminus\{\emptyset\}$ ist es nicht möglich,
    explizit eine Auswahlfunktion hinzuschreiben; das Auswahlaxiom
    garantiert jedoch ihre Existenz.
    \item \emph{Schuhe und Socken (Russell):} Ist $X$ eine Menge von
    Paaren von Schuhen, so lässt sich eine Auswahlfunktion explizit
    angeben, indem aus jedem Paar der linke Schuh gewählt wird.
    Ist $X$ dagegen eine Menge ununterscheidbarer Paare von Socken,
    lässt sich keine solche Regel angeben -- hier benötigt man AC.
\end{enumerate}
\end{beispiel}

\begin{bemerkung}[ZF, ZFC und Unabhängigkeit]
Das Auswahlaxiom wird manchmal zum Zermelo-Fraenkelschen Axiomensystem
dazugezählt und manchmal nicht: Man unterscheidet zwischen $\mathrm{ZF}$
und $\mathrm{ZFC}$ ($\mathrm{ZF}+\mathrm{AC}$).
Nach Resultaten von Gödel und Cohen ist AC unabhängig von $\mathrm{ZF}$:
Angenommen, $\mathrm{ZF}$ ist widerspruchsfrei, dann ist sowohl
$\mathrm{ZFC}$ als auch $\mathrm{ZF}+\neg\mathrm{AC}$ widerspruchsfrei.
\end{bemerkung}

\subsubsection*{Ordnungen und das Lemma von Zorn}
\begin{definition}[Halbordnung und Totalordnung]
Eine Relation $\le$ auf einer Menge $X$ heißt \textbf{Halbordnung}, wenn sie
\begin{enumerate}
    \item \textbf{reflexiv} ist: $\forall x\in X:\ x\le x$;
    \item \textbf{transitiv} ist: $\forall x,y,z\in X:\ (x\le y\land y\le z)\implies x\le z$;
    \item \textbf{antisymmetrisch} ist: $\forall x,y\in X:\ (x\le y\land y\le x)\implies x=y$.
\end{enumerate}
Eine Halbordnung heißt genau dann \textbf{Totalordnung}, wenn sie total ist,
das heißt $\forall x,y\in X:\ (x\le y\lor y\le x)$.
\end{definition}

\begin{beispiel}
Für jede Menge $X$ ist die Teilmengenrelation $\subset$ auf der
Potenzmenge $\Pot(X)$ eine Halbordnung (im Allgemeinen keine Totalordnung).
\end{beispiel}

\begin{definition}[Extrema, minimale/maximale Elemente und Schranken]
Sei $\le$ eine Halbordnung auf einer Menge $X$ und $Y\subset X$ eine Teilmenge.
\begin{itemize}
    \item Ein \textbf{Minimum} (bzw.\ \textbf{Maximum}) von $Y$ ist ein Element
    $y\in Y$, sodass für alle $y'\in Y$ gilt: $y\le y'$ (bzw.\ $y'\le y$).
    \item Ein \textbf{minimales Element} (bzw.\ \textbf{maximales Element}) von $Y$
    ist ein Element $y\in Y$, sodass es kein $y'\in Y$ gibt mit $y'\ne y$ und
    $y'\le y$ (bzw.\ $y\le y'$).
    \item Eine \textbf{untere} (bzw.\ \textbf{obere}) \textbf{Schranke} von $Y$
    ist ein Element $x\in X$, sodass für alle $y\in Y$ gilt:
    $x\le y$ (bzw.\ $y\le x$).
\end{itemize}
\end{definition}

\begin{bemerkung}
Ist $y\in Y$ ein Maximum von $Y$, dann ist $y$ sowohl ein maximales Element
von $Y$ als auch eine obere Schranke von $Y$.
Wenn ein Maximum (bzw.\ Minimum) von $Y$ existiert, ist es eindeutig:
Sind nämlich $y_1,y_2\in Y$ Maxima, so gilt $y_1\le y_2$ und $y_2\le y_1$,
also $y_1=y_2$ wegen der Antisymmetrie.
\end{bemerkung}

\begin{beispiel}[Obere Schranken und maximale Elemente]
\begin{enumerate}
    \item Auf $X=\Q$ mit der üblichen Ordnung $\le$ hat das offene Intervall
    $Y=(0,1)=\{x\in\Q\mid 0<x<1\}$ die Zahl $1$ als obere Schranke, aber
    $1\notin Y$ ist kein Maximum von $Y$. Auch $2$ ist eine obere Schranke;
    obere Schranken sind im Allgemeinen nicht eindeutig.
    \item Auf $Y=\{\{1,2\},\{1\},\{3\},\{1,2,4\}\}\subset\Pot(\N)$ mit der
    Teilmengenrelation $\subset$ gelten
    $\{1\}\subset\{1,2\}\subset\{1,2,4\}$ und $\{1\}\subset\{1,2,4\}$.
    Daher sind sowohl $\{1,2,4\}$ als auch $\{3\}$ maximale Elemente von $Y$
    (keines ist echt in einer anderen Menge aus $Y$ enthalten). Ein Maximum
    existiert nicht, da $\{3\}$ und $\{1,2,4\}$ unvergleichbar sind.
    Die Teilmenge $\{\{1\},\{1,2\},\{1,2,4\}\}\subset Y$ ist eine
    totalgeordnete Teilmenge von $Y$.
\end{enumerate}
\end{beispiel}

\begin{satz}[Lemma von Zorn]
Eine halbgeordnete Menge $(X,\le)$, in der jede totalgeordnete Teilmenge
(das heißt je zwei Elemente in der Teilmenge sind vergleichbar) eine
obere Schranke in $X$ besitzt, besitzt ein maximales Element.
\end{satz}

\begin{bemerkung}[Äquivalenz zu AC]
Im Rahmen von $\mathrm{ZF}$ sind das Auswahlaxiom und das Lemma von Zorn
äquivalent: $\mathrm{AC}\iff\text{Zorn}$ (ohne Beweis).
(Da auch die leere Teilmenge $\emptyset\subset X$ totalgeordnet ist,
impliziert die Voraussetzung des Lemmas von Zorn insbesondere $X\ne\emptyset$.)
\end{bemerkung}
